% Shared preamble for the supporting-information documents (S1-S5).
\documentclass[10pt,letterpaper]{article}
\usepackage[margin=1in]{geometry}
\usepackage{amsmath,amssymb}
\usepackage{booktabs,array,longtable}
\usepackage[table]{xcolor}
\usepackage{hyperref}
\usepackage{microtype}
\usepackage{listings}
\usepackage{graphicx}
\usepackage[labelfont=bf,labelsep=period,justification=raggedright,singlelinecheck=off]{caption}
\definecolor{linkblue}{HTML}{1C5CAB}
\definecolor{codebg}{HTML}{F4F4F1}
\hypersetup{colorlinks=true, linkcolor=linkblue, citecolor=linkblue, urlcolor=linkblue}
\lstset{basicstyle=\ttfamily\footnotesize, backgroundcolor=\color{codebg},
  frame=none, breaklines=true, columns=fullflexible, keepspaces=true,
  xleftmargin=4pt, xrightmargin=4pt, aboveskip=6pt, belowskip=6pt}
\setlength{\parindent}{0pt}
\setlength{\parskip}{5pt}
\raggedright
\renewcommand{\tablename}{Table}
\newcommand{\sititle}[2]{{\Large\bfseries #1}\par\vspace{3pt}{\large #2}\par\vspace{2pt}
  {\small\color{black!60}Supporting information for ``Environmental fluctuation limits gene regulation unless it can be tracked''}\par\vspace{10pt}}

\begin{document}
\sititle{S2 Text.}{R-Aevol build, configuration and the yoked-control derivation}

\section*{1. Software and build}
We used the regulation-enabled build of Aevol (R-Aevol) from the \texttt{aevol-7}
branch of \url{https://gitlab.inria.fr/aevol/aevol.git}, commit \texttt{33277374}.
The default branch (\texttt{aevol-9}) does not include regulation. Regulation is a
compile-time option that defines \texttt{\_\_REGUL}; its sources are in
\texttt{src/libaevol/raevol/}. The build used CMake with
\begin{lstlisting}
cmake .. -Dwith-raevol=ON -Dwithout-stats=OFF -Dwith-x=OFF \
         -DCMAKE_BUILD_TYPE=Release -DCMAKE_POLICY_VERSION_MINIMUM=3.5 \
         -DCMAKE_CXX_FLAGS="-Xpreprocessor -fopenmp -I/opt/homebrew/opt/libomp/include" \
         -DCMAKE_EXE_LINKER_FLAGS="-L/opt/homebrew/opt/libomp/lib -lomp"
\end{lstlisting}
\texttt{-Dwithout-stats=OFF} is required for the per-generation statistics file
used in all analyses; the option is on by default and suppresses it. Building on
Apple Silicon (macOS, Apple clang 21, Homebrew \texttt{cmake}, \texttt{libomp} and \texttt{boost}) needed changes that do not affect the model. The policy flag above is required because the bundled GoogleTest declares a CMake minimum below 3.5. The OpenMP include directory had to be passed through \texttt{CMAKE\_CXX\_FLAGS}, since it did not reach every target otherwise. Line 118 of \texttt{CMakeLists.txt} hard-codes the Intel Homebrew path \texttt{/usr/local/lib/libomp.dylib}, which we changed to \texttt{/opt/homebrew/opt/libomp/lib/libomp.dylib}.
All simulations ran single-threaded, ten or twelve at a time, on one 12-core
machine.

\section*{2. Parameter file}
Each run was created with \texttt{aevol\_create} from a parameter file generated
by \texttt{replication/raevol/gen\_params\_v2.py}. The exploratory runs used \texttt{gen\_params\_v1.py} and the grain experiment \texttt{gen\_params\_grain.py}, which imports the same template. The template, with the fields that vary between runs
shown in braces, is:
\begin{lstlisting}
STRAIN_NAME             {name}
SEED                    {seed}
INIT_POP_SIZE           {1024 or 512}
WORLD_SIZE              32 {32 or 16}
INIT_METHOD             ONE_GOOD_GENE CLONE
CHROMOSOME_INITIAL_LENGTH   5000
FUZZY_FLAVOR            3
SELECTION_SCHEME        fitness_proportionate   750
POINT_MUTATION_RATE     5e-5
SMALL_INSERTION_RATE    5e-5
SMALL_DELETION_RATE     5e-5
MAX_INDEL_SIZE          6
WITH_ALIGNMENTS         false
DUPLICATION_RATE        5e-5
DELETION_RATE           5e-5
TRANSLOCATION_RATE      5e-5
INVERSION_RATE          5e-5
BACKUP_STEP             10000
RECORD_TREE             false
HILL_SHAPE_N                   4
HILL_SHAPE_THETA               0.5
DEGRADATION_RATE               1
NB_DEGRADATION_STEP            10
NB_INDIV_AGE                   20
INDIVIDUAL_EVALUATION_AGES     1 2 3 4 5 6 7 8 9 10 11 12 13 14 15 16 17 18 19 20
BINDING_ZEROS_PERCENTAGE       0.75
RANDOM_BINDING_MATRIX          false
WITH_HEREDITY                  false
PROTEIN_PRESENCE_LIMIT         1e-2
ENV_SAMPLING            300
MAX_TRIANGLE_WIDTH      0.033333333
ENV_AXIS_FEATURES       METABOLISM
ENV_VARIATION           switch_in_a_list {q, or 1.5q for yoked arms}
NB_ENVIRONMENTS {2 or 4}
\end{lstlisting}
followed by the target and signal definitions of Section 3. The exploratory runs
used \texttt{BACKUP\_STEP} and \texttt{TREE\_STEP} 10{,}000 as well, and an identical
template otherwise. With \texttt{RANDOM\_BINDING\_MATRIX false} the platform reads
\texttt{binding\_matrix.rae} from the run directory; every run used the file
distributed with \texttt{examples/dynamic\_env}, unchanged. The mutation and
rearrangement rates are those of \texttt{examples/complex\_env}, ten times those of
\texttt{examples/dynamic\_env}.

\section*{3. Targets and signals}
Each environment is a target function defined by three Gaussians
(\texttt{ENV\_ADD\_GAUSSIAN} \emph{env height mean width}):
\begin{center}
\begin{tabular}{@{}lccc@{}}
\toprule
& Gaussian 1 & Gaussian 2 & Gaussian 3\\
\midrule
Target A (height, mean, width) & 1.2, 0.52, 0.12 & $-1.4$, 0.50, 0.07 & 0.3, 0.80, 0.03\\
Target B (height, mean, width) & 1.2, 0.48, 0.12 & $-1.4$, 0.50, 0.07 & 0.3, 0.20, 0.03\\
\bottomrule
\end{tabular}
\end{center}
B is A reflected about 0.5. At creation the platform reports a geometric area of 0.152468 for each. Frozen arms define two copies of A (odd seeds) or of B (even
seeds). The two signal proteins are taken from \texttt{examples/complex\_env}
(its first two \texttt{CREATE\_SIGNAL} lines):
\begin{lstlisting}
CREATE_SIGNAL h0 w0 h1 m1 w0 h1 m0 h0 h1 w0 h0 m1 h1 h1 m1 m0 h0 w0 h1 m1 w0 h1 m0 h0 h1 w0 h0 m1 h1 h1 m1 m0 h1 m0 m1
CREATE_SIGNAL m0 h0 m1 h1 m1 w0 m0 m1 m0 h0 m1 h1 w0 h0 h0 h1 m1 m0 h1 w0 h1 h0 m1 h1 m0 w0 w0 m0 m1 w0 w0 h1 h0 w0 h1 h0 h0 m0 h0 w0 h0 m1 m0 w0 h1 w0 w0 h1 m0
\end{lstlisting}
They are attached to environments with \texttt{ENV\_ADD\_SIGNAL} \emph{env signal}:
\begin{center}
\begin{tabular}{@{}lll@{}}
\toprule
Arm & Environments & \texttt{ENV\_ADD\_SIGNAL} lines\\
\midrule
Informative & (A), (B) & 1 1; 2 2\\
Yoked & (A), (A), (B), (B) & 1 1; 2 2; 3 1; 4 2\\
No cue & (A), (B) & none\\
\bottomrule
\end{tabular}
\end{center}
Before the first batch was launched, the platform's creation log was checked for
each type of arm and reported the expected signal for each environment.

\section*{4. How environments switch}
The simulations run on the platform's SIMD engine. Its environment handler
(\texttt{SIMD\_PhenotypicTargetHandler\_R::ApplyVariation}) builds, each
generation, the sequence of environments for the 20 ages of a lifetime, starting
from the environment of the last age of the previous generation:
\begin{lstlisting}
id_new_env = list_of_old_target_id[nb_indiv_age_ - 1];
id_old_env = list_of_old_target_id[nb_indiv_age_ - 1];
for (int16_t i = 0; i < nb_indiv_age_ ; i++) {
  double env_chang = var_prng_->random();
  if (env_chang < env_switch_probability_) {
    while (id_new_env == id_old_env)
      id_new_env = var_prng_->random(nb_env_);
    id_old_env = id_new_env;
  }
  list_env_id_[i] = id_new_env;
  ...
}
\end{lstlisting}
The environment is therefore a single Markov chain over ages that continues
across generations. At each age it switches with probability $q$, and a switch
moves to one of the \emph{other} environments chosen uniformly by rejection
sampling. The sequence is shared by the whole population. The non-SIMD handler
contains a branch that, when a \texttt{deterministic} flag is set and there are
two environments, gives each lifetime a fixed half-and-half split; the flag is
initialised to \texttt{false} and never set, and \texttt{check\_simd}, which would
make the SIMD engine copy the non-SIMD sequence, is never enabled by
\texttt{aevol\_run}. The branch is therefore never executed.

During development, at each age the concentrations of all signal proteins are set
to zero and those attached to the environment active at that age are set to 0.9,
before the regulatory dynamics for that age are integrated. We checked this in the
SIMD engine's network solver (\texttt{src/libaevol/7/ExpManager\_7.cpp}), which is
the code path used, as well as in the non-SIMD implementation
(\texttt{Individual\_R::EvaluateInContext}).

\section*{5. Matching switching rates in the yoked arm}
Let an environment be a pair (target, signal). In the informative arm the two
environments are (A,~s1) and (B,~s2), so every switch changes the target and the signal together. With switching probability $q$ per age, the target changes at rate
$q$ per age and the signal at rate $q$ per age.

In the yoked arm there are four environments, (A,~s1), (A,~s2), (B,~s1) and
(B,~s2). From any of them, a switch moves to one of the three others with equal
probability. Of the three, one keeps the target and changes the signal, one
changes the target and keeps the signal, and one changes the two together. A switch therefore
changes the target with probability 2/3 and the signal with probability 2/3. With switching probability $q'$, each changes at rate $2q'/3$ per age, and setting $q' = 1.5q$ gives each a rate of $q$, as in the informative arm. Without this
correction the yoked arm would fluctuate a third less often than the informative
arm, and any difference between the two would be confounded with the rate of
fluctuation.

The two arms also match in the marginal frequency of each signal (one half) and
of each target (one half), and in the number of signal proteins present at any
age (one). They differ in the association between signal and target. In the informative arm the signal determines the target, and in the yoked arm the two
are independent. They also differ in the joint dynamics of a switch, since in the
informative arm target and signal always change together and in the yoked arm
they do so in one switch out of three. This is a direct consequence of removing
the association and cannot be matched as well.

\section*{6. Measures}
R-Aevol writes \texttt{stats/stats\_simd\_best.csv}, one row per generation, for
the individual with the highest fitness. The regulatory statistics are computed in
\texttt{src/libaevol/7/Stats\_7.cpp}:
\begin{lstlisting}
for (auto affinity: rna->affinity_list) {       // coding RNAs only
  if (affinity.enhancer_factor > 0.0) { nb_activators++; mean_activator_activity += affinity.enhancer_factor; }
  if (affinity.operator_factor > 0.0) { nb_operators++;  mean_operator_activity  += affinity.operator_factor; }
}
nb_influences_       = nb_operating_influences_ + nb_enhancing_influences_;
av_value_influences_ = (mean_activator_activity + mean_operator_activity) / (nb_activators + nb_operators);
\end{lstlisting}
The column \texttt{nb\_link} is the number of positive enhancer factors plus the number of positive operator factors over all protein--RNA affinities (an affinity with a positive factor of each kind is counted twice), and \texttt{avg\_influence} is the mean of those factors,
so the total regulatory strength is $S = \texttt{nb\_link}\times\texttt{avg\_influence}$. When
there are no links, \texttt{avg\_influence} is undefined and we set $S=0$ (an
earlier, exploratory script dropped such generations instead; the difference is
immaterial for the data reported, since they are rare). Genome size is
\texttt{nb\_functional\_genes}. Each run is summarised by the mean of each measure
over generations 25{,}001--30{,}000, and the run is the unit of analysis.

\section*{7. Runs and seeds}
\begin{center}
\small
\begin{tabular}{@{}p{0.2\linewidth}p{0.42\linewidth}p{0.3\linewidth}@{}}
\toprule
Experiment & Cells & Seeds\\
\midrule
Exploratory (72 runs) & 6 arms at $q=0.10$, population 1024 & $100000 + 1000\,a + s$, $a=0..5$, $s=1..12$\\
Pre-registered sweep, primary (168) & 4 arms at $q=0.05, 0.10, 0.20$; frozen no-cue and frozen yoked at $q=0.10$; population 1024 & $200000 + 1000\,c + s$\\
Pre-registered sweep, secondary (48) & fluctuating informative, yoked and no-cue, frozen informative; $q=0.10$; population 512 & $200000 + 1000\,c + s$\\
Grain (60) & fluctuating no-cue at $q=0.1, 0.01, 0.001, 0.0002$; frozen no-cue & $300000 + 1000\,c + s$\\
\bottomrule
\end{tabular}
\end{center}
The indices $a$ and $c$ number the arm or cell in the order written by the generator
script and $s$ the replicate. Runs were launched seed-major, so that stopping a
batch early would still leave a balanced design. Every run reached 30{,}000
generations with exit status 0.

\section*{8. Pre-registration record}
\begin{center}
\small
\begin{tabular}{@{}p{0.23\linewidth}p{0.14\linewidth}p{0.55\linewidth}@{}}
\toprule
Item & Commit & Content\\
\midrule
Sweep P1--P4 & \texttt{f985ccd} & design, predictions, decision rules; \texttt{gen\_params\_v2.py}, \texttt{analyse\_v2.py}\\
EvoSim confirmation P5 & \texttt{4bc16c3} & viability-pilot rule, analysis script\\
P5 outcome & \texttt{56fe3a0} & no grid met the rule; not run\\
Grain G1--G4 & \texttt{fa00777} & \texttt{gen\_params\_grain.py}, \texttt{analyse\_grain.py}\\
\bottomrule
\end{tabular}
\end{center}
Before each analysis ran, the SHA-256 checksum of the analysis script on the
simulation machine was compared with the committed version (\texttt{analyse\_v2.py}
prefix \texttt{d0609140a970de48}; \texttt{analyse\_grain.py} prefix
\texttt{89a8613289731999}). Each analysis was run once, on complete data. The
simulation machine's idle-sleep setting suspended runs during the first night of
the sweep; this lengthened wall time but cannot change results, which are
determined by the seed. From that point runs were protected by
\texttt{caffeinate}. No other deviation occurred.

\end{document}
